\label{sec:3.5}


Our DialogGSR is trained in a multi-stage process.
The training process consists of: (1) self-supervision through knowledge graph reconstruction, (2) training the generative subgraph retriever, and (3) optimizing the response generation model. 
These stages work in synergy to ensure the model effectively retrieves knowledge from graphs and generates coherent, knowledge-grounded responses.
We also train the response generator by minimizing response generation loss.

\paragraph{Knowledge graph reconstruction.}
Inspired by masked language modeling techniques~\cite{roberts2019exploring,devlin2018bert}, we propose a self-supervised learning approach to learn the special tokens by masking either an entity token or a relation token in the token sequence of each knowledge path and reconstructing it.
Specifically, we first sample $k$-hop path $\KG^\prime$ from the knowledge source graph $\KG$ and convert it into token sequence~$\boldsymbol{z}_{\KG^\prime}$.
During training, we randomly mask out either an entity token or a relation token from the sequence. 
The loss is formulated as
\begin{equation}
    \mathcal{L}_{\text{GraphRecon}} = -\log p(\zb_{\KG^\prime} | \hat{\zb}_{\KG^\prime}),
\end{equation}
where $\boldsymbol{z}_{\KG^\prime}$ is the token sequence of a sampled path and $\hat{\boldsymbol{z}}_{\KG^\prime}$ is its randomly masked sequence.
For example, a knowledge triplet $\zb_p =\langle$ `Scarlet Letter', `written by', `N.Hawthorne' $\rangle$ can be randomly masked as 
\begin{equation*}
\begin{split}
&\langle \texttt{<M>, `written by', `N.Hawthorne'} \rangle \\
&\langle \texttt{`Scarlet Letter', <M>, `N.Hawthorne'} \rangle\\
&\langle \texttt{`Scarlet Letter', `written by', <M>} \rangle.
\end{split}
\end{equation*}
Note that masking is done at the entity or relation level as done in \citet{roberts2019exploring}. 
By minimizing the graph reconstruction loss, our framework self-supervise the special tokens $\texttt{[Head],[Int],[Rev],[Tail]}$ in \eqref{eq:kpg}, resulting in better knowledge graph representations.
All the other parameters are frozen during this stage.

\paragraph{Knowledge subgraph retrieval.}
We train our generative subgraph retriever (GSR) to identify relevant subgraphs for dialog generation.
Unlike conventional retrieval methods, our approach frames retrieval as a generation task, enabling a more seamless integration with the dialogue context. The loss is defined as follows:
\begin{equation}
    \mathcal{L}_{\text{Ret}} = \mathbb{E}_{\TokenSeq} \left [ -\log p\left({\KG}^\star | \TokenSeq \right) \right ] 
\end{equation}
where $\KG^{\star}$ is the gold subgraph and $\TokenSeq$ is the dialog context
We use cross-entropy loss to train the retriever, ensuring it generates subgraphs that are both relevant and informative.
\paragraph{Response generation.}
The final stage of training DialogGSR is response generation.
We generate dialog responses with dialog history $\TokenSeq$ and context-relevant knowledge subgraphs $\hat\KG$ retrieved from GSR. 
The response generation loss is defined as follows:
\begin{equation}
\mathcal{L}_{\text{Gen}} = \mathbb{E}_{\xb} \left [ -\log p\left(\OutSeq^\star | \TokenSeq, \hat{\KG} \right) \right ],
\end{equation}
where $\OutSeq^\star$ is the golden response.
